\section{Deferred proofs for the resolution side}\label{app:original}

This appendix proves five lemmas of Section~\ref{sec:resolution}: the regularity of the family $V_s$ (Lemma~\ref{lem:Ufan}(3)),
the structure of the regions (Lemma~\ref{lem:regions}), the positive section (Lemma~\ref{lem:pos}), the local fundamental
groups (Lemma~\ref{lem:pi1}) and the admissible cover (Lemma~\ref{lem:AC}). The notation is that of
Section~\ref{sec:resolution}; in particular $\check h_m=\check h(m)$, which is odd and positive for $m\neq0$.

\subsection{Regularity of the family: proof of Lemma~\ref{lem:Ufan}(3)}\label{app:Ufan}
\emph{Active sets.} For $C_0$ small, at every point the active monomials are vertices of one simplex of the triangulation of $\Delta^\circ$ induced by the heights, the cone over $\mathcal T$ from $0$, or, on an orbit $O(\sigma)$, of its restriction to the face $\sigma^*$. Indeed, the sets of monomials attaining $\max_iT_i$ at some point are exactly the cells of that triangulation; if a set $S$ of monomials is not contained in a cell, then $\max_iT_i-\min_{m\in S}T_m$ is a convex piecewise linear function on $N_\R$ without zeros, hence bounded below by a positive constant, and we take $C_0$ smaller than these finitely many constants. Since $\mathcal T$ is unimodular (Theorem~\ref{thm:U}), each such simplex is unimodular, so for any vertex $m_0$ of it the differences $m-m_0$ with its other vertices $m$ are part of a basis of $M$.

\emph{The computation.} Fix $s\in[0,1]$, a cone $\sigma$ with $O(\sigma)$ met by $V_s$, and $x\in V_s\cap O(\sigma)$. In Cox coordinates the defining function restricted to $O(\sigma)$ is $\sum_{m\in\sigma^*}\beta^{(s)}_m(-R)^{-\check h_m}\prod_{\varrho\notin\sigma(1)}Z_\varrho^{\langle m,n_\varrho\rangle+1}$ with $\beta^{(s)}_m=(1-s)+s\beta_m$, and on $O(\sigma)$ the factors $\beta(T_i-T_m)$ of $\beta_m$ with $i\notin\sigma^*$ equal $1$, by the argument of Lemma~\ref{lem:Ufan}(1). Let $A\subseteq\sigma^*$ be the active set at $x$, and let $m_1\in\sigma^*$ maximise $T_m(x)$; then $T_i-T_{m_1}<C_1$ near $x$ for all $i$, so $\beta_{m_1}=1$ and $\beta^{(s)}_{m_1}=1$ near $x$. A monomial $m\notin A$ has $\beta_m=0$ near $x$, hence $T_m\le T_{m_1}-C_0$ there. If $A=\{m_1\}$, the defining function is the $m_1$-term times $1+O(R^{-C_0})$, which does not vanish; so there is $m_0\in A\setminus\{m_1\}$. By the first paragraph the differences $m-m_0$, $m\in A\setminus\{m_0\}$, are part of a basis of $M$; they lie in the saturated sublattice $\sigma^\perp$, hence are part of a basis of it, and there is a cocharacter $\nu$ of the torus of $O(\sigma)$ with $\langle m_1-m_0,\nu\rangle=1$ and $\langle m-m_0,\nu\rangle=0$ for $m\in A\setminus\{m_0,m_1\}$. Let $\partial_{\rm rad}$ and $\partial_{\rm ang}$ be its radial and angular fundamental vector fields, so that $\partial_{\rm rad}x^m=\langle m,\nu\rangle x^m$, $\partial_{\rm ang} x^m=i\langle m,\nu\rangle x^m$, $\partial_{\rm rad}T_m=\langle m,\nu\rangle/\log R$ and $\partial_{\rm ang} T_m=0$.

Put $\tilde x_m=(-R)^{\check h_{m_0}-\check h_m}x^{m-m_0}$, a function on $O(\sigma)$ with $|\tilde x_m|=R^{T_m-T_{m_0}}$, and divide the defining function by the $m_0$-term:
\[
G=\sum_{m\in\sigma^*}\beta^{(s)}_m\tilde x_m=\beta^{(s)}_{m_0}+\tilde x_{m_1}+\sum_{m\in A\setminus\{m_0,m_1\}}\beta^{(s)}_m\tilde x_m+E,\qquad E=(1-s)\sum_{m\in\sigma^*\setminus A}\tilde x_m .
\]
Let $r=|\tilde x_{m_1}|\ge1$. Then $|\tilde x_m|\le r$ for all $m$, and $|\tilde x_m|\le R^{-C_0}r$ for $m\notin A$. The cut-off factors depend only on moduli, so $\partial_{\rm ang}\beta^{(s)}_m=0$; each $\beta_m$ is a product of $|\Delta^\circ\cap M|$ factors $\beta(T_i-T_m)$, and $\partial_{\rm rad}(T_i-T_m)=\langle i-m,\nu\rangle/\log R$, so $|\partial_{\rm rad}\beta^{(s)}_m|\le C_\beta/\log R$ with $C_\beta$ depending only on $\beta$, $\Delta^\circ$ and $\nu$; for $m\notin A$, $\beta^{(s)}_m=1-s$ is constant near $x$. The functions $\tilde x_m$, $m\in A\setminus\{m_1\}$, are annihilated by $\partial_{\rm rad}$ and $\partial_{\rm ang}$. Hence
\[
\partial_{\rm ang}G=i\,\tilde x_{m_1}+O(R^{-C_0}r),\qquad \partial_{\rm rad}G=\tilde x_{m_1}+O(r/\log R)+O(R^{-C_0}r),
\]
with constants depending only on $\beta$, $\Delta^\circ$ and the finitely many cocharacters $\nu$ that occur. The real $2\times2$ matrix $\mathbf J$ with columns $(\operatorname{Re},\operatorname{Im})$ of $\partial_{\rm rad}G$ and $\partial_{\rm ang}G$ has
\[
\det\mathbf J=\operatorname{Im}\bigl(\overline{\partial_{\rm rad}G}\,\partial_{\rm ang}G\bigr)=r^2\bigl(1+O(1/\log R)\bigr),
\]
which is positive for $R$ large, uniformly in $s$ and $x$. Since $G$ is the defining function times a nonvanishing function, at $x$ the derivatives of the defining function along $\partial_{\rm rad}$ and $\partial_{\rm ang}$ span $\C$. This is (3); it also shows that the differential of the defining function is nonzero at $x$ and onto $\C$ on $T_xO(\sigma)$, so $V_s$ is smooth at $x$ and transverse to $O(\sigma)$. In the notation of \cite[proof of Thm.~3.9]{Ya16}, with $\tilde x_1=\tilde x_{m_1}=r_1e^{i\theta_1}$ and $\partial_{r_1}=r_1^{-1}\partial_{\rm rad}$, $\partial_{\theta_1}=\partial_{\rm ang}$, this is $\det\mathbf J=|c_1|^2r_1\bigl(1+O(1/\log R)\bigr)$ with leading coefficient $c_1=1$. \qed

\subsection{The regions: proof of Lemma~\ref{lem:regions}}\label{app:regions}
Fix a face $\Phi$ with tie set $A=A_\Phi$, a fan point $n$ and a chart $(\pi_\parallel,\pi_\perp)$ as in Section~\ref{ssec:res-W}.

\emph{Product-shaped sets lie in $\widehat B$.} Let $U_\parallel$ be convex and $U_\perp\subset\partial \mathcal O_\perp$ connected, and put $\widetilde U=\{y:\pi_\parallel(y)\in U_\parallel,\ \pi_\perp(y)\in U_\perp\}$, a connected set. $\widetilde U\cap\nabla^{\check h}$ is closed in $\widetilde U$. It is also open: at a point of $\widetilde U\cap\nabla^{\check h}$ where a functional $\langle m,\cdot\rangle+\check h_m$ with $m\notin A$ vanishes, the point lies on a face not containing $\Phi$, which the distance condition excludes, so near such a point $\nabla^{\check h}$ is cut out by the functionals of $A$ alone. Since it is nonempty, $\widetilde U\subset\nabla^{\check h}$, and $\widetilde U\subset\widehat B$ because some functional of $A$ vanishes wherever $\pi_\perp\in\partial \mathcal O_\perp$. So a product-shaped set is exactly its model set $\widetilde U$.

\emph{The active set.} Let $y\in\Omega_U$. A monomial $m$ active at a point $x$ with $\operatorname{Log}x=y$ is within $C_{\rm act}C_0$ of the maximum of the $T_i$ at $x$. If $m\notin A$, the point $\varpi(y)$ would then lie within $C_{\rm Lip}(C_{\rm ch}\eta+C_{\rm act}C_0)=\ell_0$ of a cell on which $m$ attains the maximum, i.e.\ within $\ell_0$ of a face not containing $\Phi$; this contradicts the distance condition. Along the tentacles one uses in addition that the set of monomials attaining the maximum at a point of a tentacle is contained in that at its image under $\varpi$, because along the tentacle direction $n$ the terms of $A\setminus0$ grow at the same rate and all others at most at that rate. Hence on $\operatorname{Log}^{-1}(\Omega_U)$ only monomials of $A$ are active, and by Lemma~\ref{lem:Ufan}(3) the equation of $W$ there depends only on the characters $x^a$, $a\in A$: if $\pi_A$ denotes the homomorphism from $N\otimes\C^*$ onto the torus $T_A$ with characters $x^a$, $a\in A\setminus0$, then $W\cap\operatorname{Log}^{-1}(\Omega_U)=\pi_A^{-1}(L)\cap\operatorname{Log}^{-1}(\Omega_U)$, where $L\subset T_A$ is the localized hyperplane in the variables $x^a$.

\emph{The product.} Write $N=A^\perp\oplus\Z n\oplus C$ and $N\otimes\C^*=T_\parallel\times T_\perp$ accordingly, $T_\parallel=A^\perp\otimes\C^*$, $T_\perp=(\Z n\oplus C)\otimes\C^*$. The pairing identifies $\Z n\oplus C$ with the dual of $\Z^{A\setminus0}$, so $\pi_A|_{T_\perp}$ is an isomorphism onto $T_A$; it extends over the partial compactification along $n$, which is the part of the toric chart of the ray $n$ lying over $T_\parallel$. The logarithm splits accordingly: $\pi_\parallel(\operatorname{Log}x)$ depends only on the $T_\parallel$-component $k$, and $\pi_\perp(\operatorname{Log}x)$ only on the $T_\perp$-component, hence only on $\ell=\pi_A(x)$. Therefore
\[
W_U=\{(k,\ell)\in T_\parallel\times\bar L:\ -\log|k|/\log R\in U_\parallel^\eta,\ \pi_\perp(\ell)\in\Omega_\perp\}=(A^\perp\otimes U(1))\times U_\parallel^\eta\times L_U,
\]
with $L_U=\{\ell\in\bar L:\pi_\perp(\ell)\in\Omega_\perp\}$, $\bar L$ including the limit points on $D_n$.

\emph{The capped localized hyperplane and its truncation.}
Write $z_0=1$ and $z_i=(-R)^{-\check h_{a_i}}x^{a_i}$ for
$A\setminus0=\{a_1,\ldots,a_d\}$. Thus $p_i=-\log_R|z_i/z_0|$.
The partial compactification at $D_n$ identifies the transverse ambient space with
\[
 Z_d=\{[z_0:\cdots:z_d]\in\mathbb P^d:z_i\ne0\text{ for }1\le i\le d\},
\]
and the divisor $D_n$ is $z_0=0$. Let $Y_d\subset Z_d$ be the capped localized
hyperplane $\sum_{j=0}^d\beta_jz_j=0$, where
$\beta_j=\prod_k\beta(\log_R|z_k/z_j|)$, with its usual extension over $z_0=0$.

For completeness, $Y_d$ has the homotopy type of $(\mathbb C^*)^{d-1}$.
Apply the regularity calculation of Appendix~\ref{app:Ufan} to the standard
simplex and the family $\sum_j((1-s)+s\beta_j)z_j=0$ in $\mathbb P^d$.
At a zero on each toric orbit, two real torus vector fields have derivatives
spanning $\mathbb C$. Solving the isotopy equation with these fields and
patching by a partition of unity gives a vector field tangent to every toric
orbit; compactness of $\mathbb P^d$ gives its flow for $0\le s\le1$.
Thus the isotopy preserves every coordinate divisor and restricts to $Z_d$.
The ordinary hyperplane in $Z_d$ projects isomorphically to
$\{[z_1:\cdots:z_d]:z_i\ne0\}\cong(\mathbb C^*)^{d-1}$.
This also identifies its fundamental group with that of $Z_d$ via inclusion.

We give a deformation retraction which respects the actual logarithmic region.
For $d=2,3$, choose $C_0<a<\eta$, and set
\[
 M(z)=\max_{0\le j\le d}|z_j|,\qquad
 \delta_i(z)=\log_R\frac{M(z)}{|z_i|}\quad(1\le i\le d).
\]
The ratios defining $\delta_i$ are well defined on $Z_d$. Define
\[
 H_t[z_0:z_1:\cdots:z_d]
   =[z_0:z_1R^{t(\delta_1-a)_+}:\cdots:z_dR^{t(\delta_d-a)_+}].
\]
The maximum $M$ is unchanged. Every changed monomial has gap at least $a>C_0$
throughout this homotopy, so its cutoff coefficient vanishes; its contribution
to each active coefficient is a factor equal to $1$. The other monomials and
their coefficients do not change. Consequently $H_t$ preserves $Y_d$ exactly.
It is continuous at $z_0=0$ and is a strong deformation retraction onto
$K_a=\{z\in Y_d:\delta_i\le a\text{ for all }i\}$.

On the dense torus the logarithmic map at $t=1$ is
\[
 f_a(p)_i=\min\{p_i,\mu(p)+a\},\qquad
 H_t^{\log}(p)=(1-t)p+t f_a(p).
\]
The map $f_a$ is $1$-Lipschitz in the sup-norm: both $p_i$ and $\mu(p)+a$ are
$1$-Lipschitz, and their minimum is $1$-Lipschitz. It preserves $S_{d,r}$,
since it leaves the minimum and all coordinates attaining it unchanged and
only decreases $p_i-\mu(p)$. The same holds for $H_t^{\log}$, which is also
$1$-Lipschitz. Hence $H_t^{\log}$ preserves $S_{d,r}^{\eta}$.
For a point of $K_a$ in the dense torus, the point
$q=\mu(p)(1,\ldots,1)$ belongs to $S_{d,r}$ and satisfies
$|p-q|_\infty\le a<\eta$. The capped points follow by continuity and the local
product structure at $z_0=0$. Thus $K_a\subset L_U$, and both $Y_d$ and $L_U$
strongly deformation retract onto $K_a$. This proves
$L_U\simeq(\mathbb C^*)^{d-1}$.
For $d=1$ the localized two-term equation has its unique root at $p_1=0$,
so $L_U$ is a point.


\emph{$\pi_1$ and naturality.} $W_U$ is a product of a compact torus, a convex set and $L_U$, so it is aspherical with $\pi_1=A^\perp\oplus\pi_1(L_U)$. In type (a) this is $u^\perp\cap N$. In type (b), $\pi_1(L_U)=\Z\lambda$ with $\lambda$ the loop around the first leg, whose class generates $C$ modulo $n$; so $\pi_1=A^\perp\oplus C\cong N/\Z n$. In type (c), $\pi_1(L_U)$ is generated by the loops around the three rays $+e_k$ modulo their relation, $\cong C$, and again $\pi_1\cong N/\Z n$. Invariantly, in types (b), (c) the inclusion of $W_U$ into the toric chart of the ray $n$, whose fundamental group is $N/\Z n$, is an isomorphism on $\pi_1$; in type (a) the inclusion into the torus induces $\pi_1(W_U)=u^\perp\cap N\subset N$. For $W_{U'}\subset W_U$ both maps factor through these inclusions, and the induced map is the composite of the identifications, which is Gross's parallel transport: $u^\perp\cap N\to N/\Z n$ from type (a) to types (b), (c), and the identity of $N/\Z n$ between types (b) and (c). This proves (1) and (2).

\emph{Sheared charts.} Let $\pi^\lambda_\parallel=\pi_\parallel+\lambda\circ(\pi_\perp-\pi_\perp(0))$ with $\lambda\colon\R^{A\setminus0}\to(A^\perp)_\R$ linear and $\lambda(\pi_\perp(n)-\pi_\perp(0))=0$. Then $\pi^\lambda_\parallel(n)=0$, so tentacles are still products in the chart. In the splitting above the set defined with $\pi^\lambda_\parallel$ is $\{(k,\ell):-\log|k|/\log R+\lambda(\tau(\ell))\in U_\parallel^\eta,\ \pi_\perp(\ell)\in\Omega_\perp\}$, $\tau(\ell)=\pi_\perp(\ell)-\pi_\perp(0)$. The map $\mathrm{Sh}_\lambda(k,\ell)=(k\cdot\exp(\log R\,\lambda(\tau(\ell))),\ell)$ is a homeomorphism of $T_\parallel\times\bar L$, with inverse the same map for $-\lambda$, carrying the product set above onto this set. It extends continuously to the cap, since approaching $D_n$ the vector $\tau(\ell)$ diverges only in the direction $\pi_\perp(n)-\pi_\perp(0)=(-1,\dots,-1)$, which $\lambda$ kills. The maps $\mathrm{Sh}_{t\lambda}$ form a homotopy of maps into the toric chart of $n$, so (1) and (2) hold for the sheared chart. This proves (3). \qed

\subsection{The positive section: proof of Lemma~\ref{lem:pos}}\label{app:pos}
On the positive real torus, $\tilde f=R^{T_0}\bigl(\beta_0-\phi\bigr)$ with $\phi=\sum_{m\neq0}\beta_mR^{\tau_m}$ and $\tau_m=T_m-T_0$, because all terms other than the constant one are negative there ($\check h_m$ odd). Fix a unit vector $u$ and restrict to the ray $y=tu$, $t>0$. Then $\tau_m(tu)=-\check h_m-t\langle m,u\rangle$ is affine in $t$. Put $\bar\tau(t)=\max_{m\ne0}\tau_m(tu)$, a convex function with $\bar\tau(0)\le-1$, tending to $+\infty$ because $0$ is an interior point of $\Delta^\circ$.

(i) \emph{Signs away from the band $|\bar\tau|<C_0$.} If $\bar\tau\le-C_0$, every $m\neq0$ has $\beta(T_0-T_m)=\beta(-\tau_m)=0$, so $\beta_m=0$, and every factor of $\beta_0=\prod_i\beta(\tau_i)$ equals $1$; so $\beta_0-\phi=1$. If $\bar\tau\ge C_0$, the factor $\beta(\tau_{m^*})$ of $\beta_0$ vanishes for the maximal $m^*$, while $\beta_{m^*}=1$; so $\beta_0-\phi\le-R^{\bar\tau}<0$.

(ii) \emph{The band is an interval.} By convexity, for $t\ge t_-$, where $\bar\tau(t_-)=-C_0$, the slope of $\bar\tau$ is at least $(1-C_0)/t_{\max}$. So the band meets the ray in an interval of length at most $2C_0t_{\max}/(1-C_0)$.

(iii) \emph{Monotonicity in the band.} If $\beta_m\neq0$ then $\tau_m\ge \bar\tau-C_0\ge-2C_0$, and then $\tau_m'=-\langle m,u\rangle=(\tau_m+\check h_m)/t\ge c:=(1-2C_0)/t_{\max}>0$. Each factor $\beta(\tau_i)$ of $\beta_0$ is non-constant only where $\tau_i\in(C_1,C_0)$, where $\tau_i'>0$; so $\beta_0$ is non-increasing. Write $\beta_m=\beta(-\tau_m)\prod_{i\ne0}\beta(\tau_i-\tau_m)$. The first factor is non-decreasing along the ray, and the derivative of the product is bounded by a constant $C'_\beta$ independent of $R$; hence $\beta_m'\ge-C'_\beta\,\beta(-\tau_m)$. The maximal term has $\beta_{m^*}=\beta(-\bar\tau)$, and every active $m$ has $\tau_m\le \bar\tau$, so $\beta(-\tau_m)\le \beta(-\bar\tau)$ and $R^{\tau_m}\le R^{\bar\tau}$. Therefore
\[
\phi'=\sum_mR^{\tau_m}\bigl(\beta_m'+\beta_m\tau_m'\log R\bigr)\ \ge\ \beta(-\bar\tau)R^{\bar\tau}\bigl(c\log R-C'_\beta|\check{\mathcal A}|\bigr)>0
\]
on the band for $R$ large, since $\beta(-\bar\tau)>0$ there. So $(\beta_0-\phi)'<0$ on the band.

Hence on each ray $\tilde f$ has exactly one zero on the positive reals, and it is transversal. By the implicit function theorem the zero depends continuously (indeed smoothly) on $u$, $W_{>0}$ is a smooth hypersurface of $(\R_{>0})^4$, and $\operatorname{Log}(W_{>0})$ is a radial graph over the sphere of directions lying in the band.

(iv) \emph{Smoothness of $W$.} At a point $x\in W_{>0}$ the derivative of $\tilde f$ along the radial direction of the ray is the nonzero real number of (iii). Along the angular one-parameter group $\theta\mapsto e^{i\theta u}x$ the cut-off functions, which depend only on moduli, are constant, and
\[
\frac{d}{d\theta}\Big|_{\theta=0}\tilde f(e^{i\theta u}x)=\sum_{m\neq0}\beta_m\,i\langle m,u\rangle\,(-R)^{-\check h_m}x^m=i\sum_{m\neq0}\beta_mR^{T_m}\bigl(-\langle m,u\rangle\bigr),
\]
since the terms with $m\neq0$ are negative at $x$. By (iii) every $m$ with $\beta_m\neq0$ has $-\langle m,u\rangle\ge c>0$, and the maximal $m^*$ has $\beta_{m^*}=\beta(-\bar\tau)>0$ on the band; so this derivative is $i$ times a positive number. The two derivatives span $\C$ over $\R$, so $d\tilde f$ is onto at $x$ and $W$ is smooth along $W_{>0}$.

Since $\nabla^{\check h}$ is star-shaped from $0$ and $\bar\tau$ vanishes on each ray exactly at its point of $\partial\nabla^{\check h}$, radial projection identifies $\operatorname{Log}(W_{>0})$ with $\widehat B$. \qed

\subsection{Local fundamental groups: proof of Lemma~\ref{lem:pi1}}\label{app:collar}
Let $d=|\tau'|=\dim\tau$. We first treat the case $\dim\tau+\dim\omega=3$, in which $\dim\omega=3-d$ and $N_\omega$ has rank $4-d$, and then the case $\dim\tau=\dim\omega=1$.

\emph{Coordinates.} The characters $u-u_0$, $u\in\tau'$, are part of a basis of $M$, since $\mathcal T$ is unimodular, and they vanish on $N_\omega$, since every vertex of $\tau$ pairs to $-1$ with every lattice point of $\omega$. As $N_\omega^\perp$ is saturated, they are part of a basis of it, and a basis if its rank $3-\dim\omega$ equals $d$. Pairing with a basis of $N_\omega^\perp$ maps $N$ onto the dual lattice, with kernel $N_\omega$ (both lattices being saturated); elements $\nu_u$ (and $\nu_\kappa$) as in the lemma are obtained from preimages of the dual basis by adding multiples of a lattice point $n$ of $\omega$, which has $\langle u_0,n\rangle=-1$ and pairs to zero with $N_\omega^\perp$. Hence $N=N_\omega\oplus\bigoplus_u\Z\nu_u$ in the first case. Accordingly $X_\omega=Y_\omega\times(\C^*)^{\tau'}$, where $Y_\omega$ is the smooth toric variety of the cones of $\mathcal F'$ in $\operatorname{cone}(\omega)$, regarded in $N_\omega$, and the $\mathfrak t_u$ are the coordinates of the second factor; $\operatorname{Log}(y,\mathfrak t)=\operatorname{Log}_\omega y+\sum_u(-\log_R|\mathfrak t_u|)\nu_u$. The character $\chi=x^{-u_0}$ vanishes on the $\nu_u$, so it is a function on $Y_\omega$; it is regular and vanishes to first order along every toric divisor $D_n$ of $Y_\omega$, because $\langle -u_0,n\rangle=1$ for every lattice point $n$ of $\omega$; and $|\chi|=R^{\langle u_0,\operatorname{Log}_\omega y\rangle}$. Let $Y_{\Omega'}$ be the set of $y\in Y_\omega$ whose logarithm, modulo $\operatorname{span}(\sigma)$ on the orbit of a cone $\sigma\subseteq \operatorname{cone}(\omega)$, lies in the image of $\Omega'$; it is open, because $\Omega'$ is open and $\Omega'+\operatorname{cone}(\omega)=\Omega'$. Let $T_I=\{\mathfrak t:-\log_R|\mathfrak t_u|\in I_u\}$.

\emph{The equation.} On $\operatorname{Log}^{-1}(\Omega)$ only $0$ and the vertices of $\tau$ are active, and the factors $\beta(T_i-T_m)$ of $\beta_m$ with $i$ inactive equal $1$ for active $m$, since there $T_i\le\max_jT_j-C_0<T_m$. Moreover $T_{u_0}-T_0=-\check h_{u_0}-\log_R|\chi|$ and $T_u-T_{u_0}=\check h_{u_0}-\check h_u+\log_R|\mathfrak t_u|$. Multiplying $\tilde f$ by $x^{-u_0}$ therefore gives
\[
W_\Omega=\{(y,\mathfrak t)\in Y_{\Omega'}\times T_I:\ (\chi(y),\mathfrak t)\in L\},\qquad L=\Bigl\{(\chi,\mathfrak t):\ \beta_0\chi=\alpha_{u_0}+\sum_{u\in\tau'}\alpha_u\mathfrak t_u\Bigr\},
\]
with $\alpha_{u_0}=\beta_{u_0}R^{-\check h_{u_0}}$, $\alpha_u=\beta_uR^{-\check h_u}$ nonnegative functions of $(|\chi|,|\mathfrak t|)$; at $\chi=0$, where $\beta_0=0$, the equation is $\alpha_{u_0}+\sum_u\alpha_u\mathfrak t_u=0$. Let $\Pi(y,\mathfrak t)=(\chi(y),\mathfrak t)$ and $L_\Omega=\Pi(W_\Omega)$.

\emph{The cap.} Put $r_\chi=R^{-\check h_{u_0}-C_0}$; we call $W_\Omega\cap\{|\chi|<r_\chi\}$ the \emph{cap}. Where $|\chi|<r_\chi$, $T_0-T_{u_0}<-C_0$, so $\beta_0=0$ and $\alpha_{u_0},\alpha_u$ depend only on $|\mathfrak t|$. Hence, with $L_0=\{\mathfrak t\in T_I:\alpha_{u_0}+\sum_u\alpha_u\mathfrak t_u=0\}$ and $Y_J=\{y\in Y_{\Omega'}:|\chi(y)|\in J\}$,
\[
W_\Omega\cap\{|\chi|<r_\chi\}=Y_{[0,r_\chi)}\times L_0,\qquad W_\Omega\cap\{r_\chi/2<|\chi|<r_\chi\}=Y_{(r_\chi/2,r_\chi)}\times L_0 .
\]
$L_0$ is a point if $d=1$ and a truncated localized line, a pair of pants, if $d=2$; it is connected.

\emph{The factor $Y$.} The part $Y^\circ$ of $Y_{[0,r_\chi)}$ in the dense torus is $(N_\omega\otimes U(1))\times\{\lambda\in \Omega':\langle u_0,\lambda\rangle<\log_Rr_\chi\}$, the second factor convex, so $\pi_1(Y^\circ)=N_\omega$; likewise $\pi_1(Y_{(r_\chi/2,r_\chi)})=N_\omega$. The complement $Y_{[0,r_\chi)}\setminus Y^\circ$ is the union of the smooth divisors $D_n\cap Y_{[0,r_\chi)}$, each nonempty because $\Omega'+\operatorname{cone}(\omega)=\Omega'$ and $\langle u_0,n\rangle=-1$. By general position a loop in the manifold $Y_{[0,r_\chi)}$ can be pushed into $Y^\circ$, and a null-homotopy can be made transverse to the divisors, meeting them in finitely many points of their smooth parts; so $\pi_1(Y_{[0,r_\chi)})$ is the quotient of $N_\omega$ by the classes of the meridians of the $D_n$, which are the orbits of the circles $\Z n\otimes U(1)$. The rays $n$ of a maximal cone of $\mathcal F'$ in $\operatorname{cone}(\omega)$ form a basis of $N_\omega$, since that cone is unimodular (Lemma~\ref{lem:Ufan}(2)). Hence $Y_{[0,r_\chi)}$ is simply connected.

\emph{The part away from the cap.} Fix a ray $n_0$ of $Y_\omega$; then $\langle -u_0,n_0\rangle=1$ and $N_\omega=N^0_\omega\oplus\Z n_0$ with $N^0_\omega=N_\omega\cap u_0^\perp$. Writing points of $Y^\circ$ as $(\theta,\lambda)\in(N_\omega\otimes U(1))\times \Omega'$ and $\theta=(\theta_{N^0_\omega},\theta_0)$, the phase $\theta_0$ is determined by $\arg\chi$, so
\begin{gather*}
W_{>r_\chi/2}:=W_\Omega\cap\{|\chi|>r_\chi/2\}\cong(N^0_\omega\otimes U(1))\times \mathcal S,\\
\mathcal S=\{(\lambda,\ell)\in \Omega'\times L_{>r_\chi/2}:\ \langle u_0,\lambda\rangle=\log_R|\chi(\ell)|\},
\end{gather*}
$L_{>r_\chi/2}=L_\Omega\cap\{|\chi|>r_\chi/2\}$. The projection $\mathcal S\to L_{>r_\chi/2}$ has nonempty convex fibres $\Omega'\cap\{\langle u_0,\cdot\rangle=\log_R|\chi|\}$, and it has a continuous section $\ell\mapsto\lambda(\log_R|\chi(\ell)|)$: choose points $p_j\in \Omega'$, $j\ge0$, with $\langle u_0,p_j\rangle$ increasing to $\sup_{D'}\langle u_0,\cdot\rangle$, let $\lambda(c)$ run linearly from $p_j$ to $p_{j+1}$ for $c$ between their values, and put $\lambda(c)=p_0+(\langle u_0,p_0\rangle-c)n_0$ below, which lies in $p_0+\operatorname{cone}(\omega)\subseteq \Omega'$. The straight-line homotopy in the fibres retracts $\mathcal S$ onto the section, so $\pi_1(W_{>r_\chi/2})=N^0_\omega\times\pi_1(L_{>r_\chi/2})$.

\emph{Van Kampen.} $W_\Omega$ is the union of the open sets $W_{[0,r_\chi)}=Y_{[0,r_\chi)}\times L_0$ and $W_{>r_\chi/2}$, with connected intersection $Y_{(r_\chi/2,r_\chi)}\times L_0$, whose fundamental group $N_\omega\times\pi_1(L_0)$ maps to $\pi_1(L_0)$ by the projection, and to $N^0_\omega\times\pi_1(L_{>r_\chi/2})$ by the identity on $N^0_\omega$, by $n_0\mapsto\gamma_\chi$, the loop in $L_{>r_\chi/2}$ that turns $\arg\chi$ once with $\mathfrak t$ fixed in $L_0$, and by the inclusion of $\{\chi\}\times L_0$ on $\pi_1(L_0)$. So
\[
\pi_1(W_\Omega)\cong\pi_1(L_{>r_\chi/2})/\langle\!\langle\gamma_\chi\rangle\!\rangle\cong\pi_1(L_\Omega),
\]
the second isomorphism by van Kampen for $L_\Omega=L_{>r_\chi/2}\cup(\{|\chi|<r_\chi\}\times L_0)$; both are induced by $\Pi$. The same decompositions show that $W_\Omega$ is path connected, and it is open in $W$ because $Y_{\Omega'}$ and $T_I$ are open.

\emph{The interval regions.}
In the notation of Appendix~\ref{app:collar}, enumerate $\tau'=\{u_2,\ldots,u_{d+1}\}$ and use projective monomial
coordinates
\[
 z_0=\chi,\qquad z_1=R^{-\check h_{u_0}},\qquad
 z_i=R^{-\check h_{u_i}}\mathfrak t_{u_i}\quad (i=2,\ldots,d+1).
\]
The signs can equivalently be absorbed in these coordinates; they do not affect
any modulus below. The image region defining $L_\Omega$ is described by
\[
 c=\log_R|z_0/z_1|<b_0,\qquad
 b_i=\log_R|z_i/z_1|\in J_i\quad (i=2,\ldots,d+1),
\]
where $b_0\ge4C_0$ (possibly $+\infty$) and every $J_i$ contains
$[-4C_0,4C_0]$. The equality of this image region with these intervals follows
from the nonempty convex fibres of the projection from $\Omega'$;
$\Omega'+\operatorname{cone}(\omega)=\Omega'$ implies that the $c$-interval
has no lower endpoint. Apply the same homotopy $H_t$ with $a=2C_0$, clamping
all the nonconstant monomials $z_1,\ldots,z_{d+1}$ and leaving $z_0$ unchanged.
Writing $h_i=(\delta_i-a)_+$, its action is
\[
 b_i(t)=b_i+t(h_i-h_1),\qquad c(t)=c-th_1.
\]
Since $\delta_i=\delta_1-b_i$, each $b_i(t)$ lies between $b_i$ and $0$;
hence it remains in $J_i$. Also $c(t)\le c$, so the upper bound is preserved.
At $t=1$, all $|b_i|\le a$ and $c\le a$, so the clamped core is contained in
the image region. Therefore the full capped localized hyperplane and
$L_\Omega$ both retract onto this core. In particular
$L_\Omega\simeq(\mathbb C^*)^d$. The characters $\mathfrak t_u$ do not vanish on $L_\Omega$ and define a homomorphism $\pi_1(L_\Omega)\to\Z^{\tau'}$. It is onto. In the coordinates $\log_R|\chi|$ and $\log_R|\mathfrak t_u|$ the region above is a product of intervals, and each contains the value at the vertex of the tropical hyperplane with a margin of at least $4C_0$, because the projection of the $4C_0$-ball in $\Omega$ under a nonzero lattice covector has half-width at least $4C_0$. For $u\in\tau'$, decreasing $\log_R|\mathfrak t_u|$ from its value at the vertex by $2C_0$, with the other coordinates at their values at the vertex, gives a point of the region at which the term of $u$ lies $2C_0$ below the others, which tie. Near that point only the other monomials are active, so $L$ is there the product of the localized hyperplane of the other monomials in the remaining variables (a point if $d=1$) and an annulus in $\mathfrak t_u$; the loop that turns $\arg \mathfrak t_u$ once, with the other coordinates fixed, lies in $L_\Omega$ and maps to the basis vector of $u$. A surjective homomorphism $\Z^d\to\Z^d$ is an isomorphism. By the previous step the characters $\mathfrak t_u$ induce an isomorphism on $\pi_1(W_\Omega)$ as well.

\emph{The case $\dim\tau=\dim\omega=1$.} Here $d=1$, $\tau'=\{u\}$, and $N_\omega^\perp$ has the basis $u-u_0,\kappa$. As above $N=N_\omega\oplus\Z\nu_u\oplus\Z\nu_\kappa$, $X_\omega=Y_\omega\times\C^*_{\mathfrak t_u}\times\C^*_w$ with $w=x^\kappa$, and $\operatorname{Log}(y,\mathfrak t_u,w)=\operatorname{Log}_\omega y-\log_R|\mathfrak t_u|\,\nu_u-\log_R|w|\,\nu_\kappa$. The character $\chi=x^{-u_0}$ vanishes on $\nu_u$ and $\nu_\kappa$, and the functions $T_0-T_{u_0}$ and $T_u-T_{u_0}$ depend only on $|\chi|$ and $|\mathfrak t_u|$; so the equation of $W$ on $\operatorname{Log}^{-1}(\Omega)$ does not involve $w$, and
\[
W_\Omega=W_\Omega^\flat\times\{w:-\log_R|w|\in I_\kappa\},
\]
where $W_\Omega^\flat\subset Y_\omega\times\C^*_{\mathfrak t_u}$ is defined by the conditions above with $\Omega'\oplus I_u\nu_u$. The preceding steps apply to $W_\Omega^\flat$ word for word: $Y_\omega$ is the smooth toric surface $\C^2$ of $\operatorname{cone}(\omega)$, in which $\chi$ vanishes to first order along both coordinate axes, $N^0_\omega=N_\omega\cap u_0^\perp$ has rank one, $L_0$ is a point, and $L_\Omega$ is a capped localized line as in type (b) of Appendix~\ref{app:regions}. So $W_\Omega^\flat$ is open and path connected, with $\pi_1$ mapped isomorphically to $\Z$ by $\mathfrak t_u$, and the annulus contributes a factor $\Z$ detected by $w$. \qed

\subsection{The admissible cover: proof of Lemma~\ref{lem:AC}}\label{app:AC}
\emph{Data.} Let $d_S=\operatorname{dist}(\hat S\setminus \mathcal N_\Gamma,\widehat\Gamma)>0$, distances being Euclidean for the standard inner product in lattice coordinates on $N_\R$ (so $|m|\ge1$ for $0\neq m\in M$). For a two-face or an edge $\Phi$ and a fan point $n$ we use the \emph{canonical chart} $\pi_\Phi$, the projection onto $(A_\Phi^\perp)_\R$ along $J_\Phi(n)=\R n+((A_\Phi^\perp)_\R+\R n)^\perp$. $J_\Phi(n)$ is a complement of $(A_\Phi^\perp)_\R$ containing $n$, so $\pi_\Phi$ differs from a lattice chart by a shear killing $\pi_\perp(n)-\pi_\perp(0)$, and Lemma~\ref{lem:regions}(3) applies. These charts nest: for $\Phi'\supset\Phi$ with the same $n$, $J_{\Phi'}(n)\subset J_\Phi(n)$, so $\pi_\Phi=\mathrm{pr}\circ \pi_{\Phi'}$ for a linear projection $\mathrm{pr}$. For a facet $F=F_u$ fix one chart, the projection along $\R n_F$ with $\langle u,n_F\rangle=-1$. Let $C_{\rm ch}$ bound the distortion between the chart sup-norms and the Euclidean norm, and $C_{\rm norm}$ the norms of the maps $\mathrm{pr}$ between nested charts, over the finitely many charts in use; put $C_*=C_{\rm norm}C_{\rm ch}^2$, $\gamma_F=0$, $\gamma_G=1$ for two-faces, $\gamma_E=1+C_{\rm norm}(\gamma_G+1)$ for edges.

\emph{Cells and labels.} Let $\mathfrak T$ be a triangulation of $S\setminus\operatorname{int}\mathcal N_\Gamma$ for which $\hat j$ is linear on cells and which is adapted to the faces of $\nabla^{\check h}$. A cell $c$ has type (a), (b) or (c) according as the face $\Phi_c$ of lowest dimension met by $\hat j(c)$ is a facet, a two-face or an edge. For types (b), (c), $\hat j(c)\cap\Phi_c$ is connected and at distance $\ge d_S$ from $\widehat\Gamma$, and the region boundaries in the two-skeleton lie in $\widehat\Gamma$ (Lemma~\ref{lem:twofaces}); so it lies in one vertex region, that of a vertex we call $n_c$. If $c'\subset c$ both of type (b) or (c), then $n_{c'}=n_c$. Put $h_c=\sup_{\hat j(c)}|\pi_\perp|_\infty$ in the chart of $(\Phi_c,n_c)$, and for a radius $\lambda$
\[
V_c=\bigl\{y\in\widehat B:\ \pi_\parallel(y)\in(\operatorname{conv}\pi_\parallel\hat j(c))^{\lambda},\ \pi_\perp(y)\in\partial \mathcal O_\perp,\ |\pi_\perp(y)|_\infty<\gamma_{\Phi_c}(h_c+\lambda)\bigr\}
\]
for types (b), (c), and $V_c=\{y\in F:\ |\pi_F(y)-\pi_F(\operatorname{conv}\hat j(c))|<\lambda\}$ for type (a), where $\lambda=\lambda_{\dim c}$.

\emph{The condition (G0).} For a cell of type (b) or (c) let $V^0_c=\{y:\pi_\parallel(y)\in\operatorname{conv}\pi_\parallel\hat j(c),\ \pi_\perp(y)\in\partial \mathcal O_\perp,\ |\pi_\perp(y)|_\infty\le \gamma_{\Phi_c}h_c\}$, the set $V_c$ at $\lambda=0$; for type (a), $V^0_c=\operatorname{conv}\hat j(c)$. We require of $\mathfrak T$:
\begin{enumerate}
\item[(G0)] for every cell $c$, every functional $\ell_m=\langle m,\cdot\rangle+\check h_m$ with $m\notin A_{\Phi_c}$ is positive on $V^0_c$, and $V^0_c\cap\widehat\Gamma=\emptyset$.
\end{enumerate}
By compactness the margin $m(c)=\min\bigl(\min_m\min_{V^0_c}\ell_m/|m|,\ \operatorname{dist}(V^0_c,\widehat\Gamma)\bigr)$ is positive; note $\operatorname{dist}(y,F_m)\ge\ell_m(y)/|m|$.

\emph{Order of choices.} After $\hat S$, $\mathcal N_\Gamma$ and the generic isotopy: the charts and constants above; a triangulation $\mathfrak T$ satisfying (G0); then radii $\lambda_0<\lambda_1<\lambda_2$ such that $\lambda_k\ge2C_*\lambda_{k-1}$ and $\lambda_2<\min_cm(c)/(2C_{\rm ch}(\gamma_E+1))$; then $\eta\le\lambda_0/(2C_*)$ with $(C_{\rm norm}+1)\eta<\min_cm(c)/2$ and $\eta<\min_cm(c)/(4C_{\rm Lip}C_{\rm ch})$; then $C_0<\min\bigl(\eta/(4C_{\rm ch}t_{\max}),\ \min_cm(c)/(4C_{\rm act}C_{\rm Lip})\bigr)$, which makes the band of Lemma~\ref{lem:pos} narrower than $\eta/C_{\rm ch}$ and gives $\ell_0<\min_cm(c)/2$; finally $R$ large. Every constraint is an upper bound.

\emph{Admissibility (Lemma~\ref{lem:AC}(1)).} $V_c$ lies within $C_{\rm ch}(\gamma_{\Phi_c}+1)\lambda_2<m(c)/2$ of $V^0_c$, and $\ell_0<m(c)/2$; so $V_c$ is more than $\ell_0$ away from every face not containing $\Phi_c$ and from $\widehat\Gamma$. By the first step of Appendix~\ref{app:regions}, $V_c$ is a product-shaped subset of $\widehat B$. It contains $\hat j(c)$ by construction. For type (a) (G0) is automatic: each face of the convex facet $F$ is exposed by a functional that is non-negative on $F$, so it is positive on $\operatorname{conv}\hat j(c)$ if $\hat j(c)$ misses that face.

\emph{Monotonicity (Lemma~\ref{lem:AC}(2)).} Let $c'\subset c$; the lowest face of $c'$ contains that of $c$. Write $\lambda'=\lambda_{\dim c'}<\lambda=\lambda_{\dim c}$. First, $V_{c'}\subset V_c$ with a margin. If $\Phi_{c'}=\Phi_c$, the charts agree, the hull and $h$ shrink, and $\gamma(h_{c'}+\lambda')\le \gamma(h_c+\lambda)-\gamma(\lambda-\lambda')$. If $c'$ has type (a) and $c$ type (b) or (c), every $y\in V_{c'}$ is within $C_{\rm ch}\lambda'$ of $\operatorname{conv}\hat j(c')\subset\operatorname{conv}\hat j(c)$, so $\pi_\parallel(y)$ is within $C_{\rm ch}^2\lambda'$ of $\operatorname{conv}\pi_\parallel\hat j(c)$ and $|\pi_\perp(y)|\le h_c+C_{\rm ch}^2\lambda'$. If $c'$ has type (b) at $G$ and $c$ type (c) at an edge $E\subset G$, take $y\in V_{c'}$ and $w\in\operatorname{conv}\hat j(c')$ with $|\pi_G(y)-\pi_G(w)|<\lambda'$. By nesting $|\pi_E(y)-\pi_E(w)|\le C_{\rm norm}\lambda'$. Further $|y-w|_G\le\max(\lambda',r_{c'}+h_{c'})$, and $\pi^G_\perp$ is a sub-tuple of $\pi^E_\perp$ independent of $n$, so $h_{c'}\le h_c$ and
\[
|\pi^E_\perp(y)|\le h_c+C_{\rm norm}(r_{c'}+h_{c'})\le \gamma_Eh_c+C_{\rm norm}\gamma_G\lambda'\le r_c-\lambda .
\]
This is where the grading $\gamma_G<\gamma_E$ of the transverse radii is needed. In every case the room between $V_{c'}$ and the frontier of $V_c$ is at least $\lambda-C_*\lambda'\ge\lambda/2\ge C_*\eta$ in the chart norm of $c$.

Second, $\Omega_{V_{c'}}\subset\Omega_{V_c}$. Let $z$ be within $\eta$ of $\varpi^{-1}(V_{c'})$ in the chart norm of $c'$. In the $\pi_\parallel$-directions of the chart of $c$, $z$ is within $C_{\rm norm}\eta$ of $\pi_\parallel(V_{c'})$; in the $\pi_\perp$-directions of $G$ there is no loss, $\pi^G_\perp$ being a sub-tuple of $\pi^E_\perp$. In the remaining coordinate $\pi_{u_3}$ of $E$ the tentacle direction has $d\pi_{u_3}(n)=-1$, which absorbs the deviation along the tentacle; the transverse deviation is at most $C_{\rm norm}\eta$, and it is absorbed because $\ell_{u_3}$ is a foreign functional of $c'$, so $\ell_{u_3}\ge|u_3|(m(c')-C_{\rm ch}\lambda')>m(c')/2>(C_{\rm norm}+1)\eta$ on $V_{c'}$. The case of type (a) inside type (b) or (c) is the same, using $\pi_{u'}\ge(C_{\rm norm}+1)\eta$ on $V_{c'}$ for the other facets $u'$ through $\Phi_c$, which (G0) and the choice of $\eta$ give. By the margin of the first step, $z\in\Omega_{V_c}$.

\emph{Existence of a triangulation satisfying (G0).} Let $x_0$ be an edge crossing, on the edge $E$ of $\nabla^{\check h}$. Choose a convex ball $U_0$ around $x_0$ that misses $\widehat\Gamma$ (possible since $\widehat\Gamma\cap E$ is finite and $x_0\notin\widehat\Gamma$ by genericity), the other crossings, $\mathcal N_\Gamma$ and every face not containing $E$, and that lies in the open star of $x_0$ in a triangulation of $\widehat B$ for which $\hat j$ is simplicial and the faces are subcomplexes. Then $\hat S\cap U_0$ is a cone from $x_0$, and so is $\widehat B\cap U_0=\partial C_E\cap U_0$, $C_E$ the cone cut out by the three functionals of $A_E$. The cone $\hat S\cap U_0$ meets each of the three two-faces $G_k$ through $E$ in finitely many rays.

\emph{Foreign functionals near $x_0$.} Let $y\in\partial C_E\cap U_0$ and $m\notin A_E$. Then $\ell_m(x_0)>0$. If some foreign functional vanished on $[x_0,y]$, let $z$ be its first zero; at $z$ all other foreign functionals and those of $A_E$ are $\ge0$ ($C_E$ is convex), so $z\in\nabla^{\check h}$ lies on the facet $F_m$, a face not containing $E$, inside $U_0$, a contradiction. So on $\partial C_E\cap U_0$ only the functionals of $A_E$ can vanish, and these are linear, hence invariant under homotheties centred at $x_0$.

\emph{The first ring.} Take a polygon in $\hat S\cap U_0$ around $x_0$ with a vertex on every ray of $\hat S\cap\bigcup_kG_k$ and at every crease of the cone, so that each outer edge meets at most one two-face and each cone triangle lies in $\hat S$. The first ring consists of the triangles $[x_0,p_i,p_{i+1}]$; it may later be subdivided only by coning from $x_0$, so every cell containing $x_0$ keeps type (c). (Splitting a ring triangle at an interior point of $[x_0,p_i]$ would create a type-(b) cell touching $G_k$ near $x_0$ whose set $V^0$ can leave $\nabla^{\check h}$.) For a ring cell $c$ and $y\in V^0_c$, $|\pi_\parallel(y)-\pi_\parallel(x_0)|\le C_{\rm ch}R_c$ and $|\pi_\perp(y)|\le \gamma_Eh_c\le \gamma_EC_{\rm ch}R_c$, where $R_c=\sup_{\hat j(c)}|z-x_0|\le s$ for a ring of radius $s$. So $|y-x_0|\le C_{\rm ch}^2(\gamma_E+1)s$, and for $s$ small $V^0_c\subset\partial C_E\cap U_0$, where (G0) holds by the previous step.

\emph{Type-(b) cells near $x_0$.} Let $c'$ be a type-(b) cell at $G$ with $\hat j(c')\subset\frac12U_0$, and put $t=\operatorname{dist}(x_0,\hat j(c'))>0$. Since the link of the cone $\hat S\cap U_0$ is compact and crosses $G$ transversally, there are $\varphi_1,c_\gamma>0$ with $\ell_{u_3}(z)\ge c_\gamma|z-x_0|$ for $z\in\hat S\cap U_0$ with $|\pi^G_\perp(z)|\le\varphi_1|z-x_0|$, $u_3$ the element of $A_E\setminus A_G$. In the chart of $G$ let $w_1,w_2$ be the vectors with $\pi_\parallel=0$ and $d\pi^G_\perp=e_1,e_2$, and $\beta_w=\max_a|\langle u_3,w_a\rangle|$. A point of $V^0_{c'}$ has the form $y_G+\sigma w_a$ with $0\le\sigma\le \gamma_Gh_{c'}$ and $y_G$ in the hull of the points $z-\pi_a(z)w_a$, $z\in\hat j(c')$. If $h_{c'}\le\varphi_0t$, then $\ell_{u_3}\ge(c_\gamma-2\varphi_0\beta_w)t>0$ on $V^0_{c'}$ for $\varphi_0\le\min(\varphi_1,c_\gamma/(4\beta_w))$, chosen also so small that $V^0_{c'}\subset U_0$. Homothety about $x_0$ commutes with the chart of $G$, so $\varphi_0$ is uniform in $t$. Then $V^0_{c'}\subset\partial C_E\cap U_0$ and (G0) holds by the first-zero argument. The condition $h\le\varphi_0t$ is met by grading the outer polygon angularly near each ray and the cells along each ray geometrically.

\emph{The remainder and subdivision.} Triangulate the rest of $\hat S\setminus \mathcal N_\Gamma$ with the outer polygons of the rings as inner boundary. On this compact remainder, away from the crossings, take cells small compared with their distance to the faces they do not meet and to $\widehat\Gamma$; then $h_c$ is small too, the lowest face is well defined, and (G0) follows from the first-zero argument started at a point of $\Phi_c$. (G0) survives any adapted subdivision of the cells outside the rings: a subcell with the same lowest face and label has a smaller $V^0$; a type-(b) subcell of a type-(b) cell near $x_0$ has $t$ at least as large and $h$ at most as large; a type-(a) subcell satisfies (G0) automatically.

\emph{Support (Lemma~\ref{lem:AC}(3)).} $\hat j(c)$ meets $\Phi_c$, where $\pi_\perp=0$, so $h_c\le C_{\rm ch}\operatorname{diam}\hat j(c)$, and $V_c$ lies in the $r'$-neighbourhood of $\hat j(c)$ with $r'=O(\operatorname{mesh}(\mathfrak T)+\lambda_2)$. The mesh can be decreased by subdivision outside the rings and by shrinking the rings, and $\lambda_2$, $\eta$, $C_0$ are bounded only from above. \qed
